\documentclass{article}

\input{macros.tex}


\begin{document}

\blogtitle{The Shortest and Closest Vector Problems}
\blogsubtitle{The fundamental problems of lattice-based cryptography.}
\written{2026-09-13}
\updated{2026-09-13}


So far, we have only talked about properties and ways to describe lattices, but that doesn't inherently have anything to do with cryptography.
This connection is made through hiding information by utilizing the structure of lattices.
In other words, lattice-based cryptography relies on lattice \emph{problems}, such as the shortest vector problem (SVP).
SVP essentially asks questions surrounding finding the shortest vectors of a lattice given a basis.

Recall that there are an infinite number of bases that generate the same lattice.
In previous blogs, we have given relatively "simple" bases, but that is not always the case.
For instance, consider the basis
\[
B :=
\begin{pmatrix}
3\sqrt{3898}
&
-5382\sqrt{\frac{2}{1949}}
&
\frac{31195}{\sqrt{3898}}
&
\frac{15857}{3}\sqrt{\frac{2}{1949}}
\\
0
&
\sqrt{\frac{682378}{1949}}
&
-110727\sqrt{\frac{2}{664977361}}
&
\frac{676011}{\sqrt{1329954722}}
\\
0
&
0
&
\sqrt{\frac{64221}{682378}}
&
\frac{67240}{3}\sqrt{\frac{2}{21911498769}}
\\
0
&
0
&
0
&
\frac{1}{3\sqrt{128442}}
\end{pmatrix},
\]
a basis is borrowed from [1].
How would we interpret this?
What is the shortest vector in the lattice formed by this basis?
Since this is not immediately obvious (to a human, but computationally this is extremely weak and not secure), this is a bit more difficult than the bases we were previously using.




\section{Minimum Distance}
We need to collect some more lattice ``tools'' to analyze and understand SVP, especially what it means to have a shortest vector.
We discuss some linear algebra machinery such as the \emph{Euclidean norm}, sometimes denoted as the $\ell_2$ norm.
More specifically, for some vector $x = (x_1, \ldots, x_n) \in \RR^n$, the $\ell_2$ norm is
\[
\norm{x}_2 = (x_1^2 + \ldots + x_n^2)^{1/2},
\]
essentially the ordinary geometric length from the origin to $x$.
Sometimes other norms will be used too:
\[
\norm{x}_p = (\lvert x_1\rvert^p + \ldots + \lvert x_n \rvert^p)^{1/p}
\]
for $p \geq 1$, but this is not as common.
Assume we are using the Euclidean norm unless stated otherwise.


The \emph{minimum distance} of a lattice $L$ is defined as
\[
\lambda_1(L)
:=
\min_{v \in L \setminus \{0\}} \lVert v \rVert
=
\min_{x,y \in L}
\lVert x-y \rVert
\]
where $x \neq y$.
In non-notation terms, it is the minimum non-zero norm of a vector within a lattice, which can be generalized to distinct vectors.
Figure two shows how this works on a two-dimensional graph.

\blogfigure{/assets/svg/lattices/2/figure-1.svg}
{\textbf{Figure 1.} One vector with the minimum distance in two dimensions.}

\section{SVP}
Let's define an approximate version of SVP, which can also be used for an exact version.
The $\gamma$-approximate SVP, where $\gamma \geq 1$, is defined in three variants.
This language is largely borrowed from [2], where we (somewhat) grossly simplify.


\subsection{Decision}
Decision ($\mathsf{GapSVP}_\gamma$) asks, given a lattice basis $B$ and a real $d > 0$, to distinguish between 
\[
\lambda_1 (L(B)) \leq d \text{ and } \lambda_1 (L(B)) > \gamma \cdot d
\]
The exact version $\mathsf{GapSVP}$ is when $\gamma = 1$.
This asks us to tell whether $d$ is greater than or equal to the minimum distance, or (approximately) less than the minimum distance.
Figure 2 shows how the problem looks on a number line.

\blogfigure{/assets/svg/lattices/2/figure-2.svg}
{\textbf{Figure 2.} $\mathsf{GapSVP}_\gamma$ on a one-dimensional number line.}

\subsection{Estimation}
Estimation ($\mathsf{EstSVP}_\gamma$) asks, given a lattice basis $B$, to compute $\lambda_1 (L(B))$ up to a $\gamma$ factor, meaning in the range 
\[
[\lambda_1 (L(B)),\ \gamma \cdot \lambda_1 (L(B))].
\]
The exact version $\mathsf{EstSVP}$ asks us to calculate $\lambda_1 (L(B))$.
This is asking us to calculate the (range of the) minimum distance.

\subsection{Search}
Search ($\mathsf{SVP}_\gamma$) asks, given a lattice basis $B$, to find a nonzero $v \in L(B)$ such that 
\[
\norm{v} \leq \gamma \cdot \lambda_1 (L(B)).
\]
The exact version $\mathsf{SVP}$ asks us to calculate $v$ such that $\norm{v} = \lambda_1(L(B))$.
Essentially, it asks us to calculate the vector with the (approximate) minimum distance.


\subsection{Discussion on hardness}
While we will talk a bit more formally about hardness in a future post, it might be helpful to contextualize the differences between the problems.
For now, it suffices to say that if you can solve the search problem, you can solve the other two problems.
An intuitive way to think of this is that if you can calculate the shortest vector, you can compute $\lambda_1(L(B))$ (or the approximate version).
Thus, you can solve the decision problem.

There are a few algorithms, historic and contemporary, that are common when speaking about calculating SVP.
One example is the LLL algorithm [3], which approximates a short nonzero lattice vector.
Additionally, a randomized sieving technique [4] solves SVP in $2^{O(n)}$ time and space. 
The contemporary discrete-Gaussian-sampling algorithm [5] solves it in $2^{n+o(n)}$ time and space.
Note that the solving algorithms (not inherently LLL) are all exponential in $n$, meaning that there is no known polynomial or efficient algorithm.



\section{CVP}
The Closest Vector Problem (CVP) is closely related to SVP.
Rather than finding a shortest nonzero vector in a lattice, we are given some target $t \in \mathbb{R}^n$ and consider the distance between $t$ and the closest lattice point.
We define this distance as
\[
\dist{t,L(B)}
    =
    \min_{v \in L(B)} \norm{t-v}.
\]

The decision, estimation, and search variants of CVP are defined similarly to their respective SVP variants, replacing $\lambda_1(L(B))$ with $\dist{t,L(B)}$.
Thus, $\mathsf{GapCVP}_\gamma$ asks whether the target is within some distance $d$ of the lattice or farther than $\gamma d$.
$\mathsf{EstCVP}_\gamma$ approximates $\dist{t,L(B)}$.
$\mathsf{CVP}_\gamma$ finds a lattice vector whose distance from $t$ is within a factor $\gamma$ of the minimum possible distance, i.e., $\norm{t-v} \leq \gamma \dist{t,L(B)}$.

Intuitively, SVP asks us to find the shortest vector relative to the origin, while CVP lets us choose an arbitrary target and asks for the closest lattice point to that target.
Although they are related, they are not the same problem.
Figure 3 displays CVP graphically.
\blogfigure{/assets/svg/lattices/2/figure-3.svg}
{\textbf{Figure 3.} An example of the distance between a point $t$ and a point in $L$.}


\subsection{Discussion on hardness}
As with SVP, the search version gives us enough information to solve the calculation and decision versions.
Additionally, SVP can be solved using a CVP ``solver'' [6], so CVP is at least as hard as SVP for $\gamma \geq 1$.

Exact CVP is known to be hard [7, 8], and there is no known polynomial-time algorithm for solving it.
This does \emph{not} imply that SVP is hard.
One important exact algorithm is the MV Voronoi-cell algorithm [9], which solves CVP deterministically in $O(2^{2n})$ time and $O(2^n)$ space.






\bib{
@inproceedings{10.1007/978-3-031-30589-4_9,
author = {Bennett, Huck and Ganju, Atul and Peetathawatchai, Pura and Stephens-Davidowitz, Noah},
title = {Just How Hard Are Rotations of $\ZZ_n$? Algorithms and Cryptography with the Simplest Lattice},
year = {2023},
isbn = {978-3-031-30588-7},
publisher = {Springer-Verlag},
address = {Berlin, Heidelberg},
url = {https://doi.org/10.1007/978-3-031-30589-4_9},
doi = {10.1007/978-3-031-30589-4_9},
booktitle = {Advances in Cryptology – EUROCRYPT 2023: 42nd Annual International Conference on the Theory and Applications of Cryptographic Techniques, Lyon, France, April 23-27, 2023, Proceedings, Part V},
pages = {252–281},
numpages = {30},
location = {Lyon, France}
}}

\bib{
@misc{peikert-lattices,
    author = {Chris Peikert},
    title = {Lattices in Cryptography},
    year = {2026},
    url = {https://github.com/cpeikert/LatticesInCryptography}
}
}

\bib{
@article{LenstraLenstraLovasz1982,
  author  = {Lenstra, Arjen K. and Lenstra, Hendrik W., Jr. and Lov{\'a}sz, L{\'a}szl{\'o}},
  title   = {Factoring Polynomials with Rational Coefficients},
  journal = {Mathematische Annalen},
  volume  = {261},
  number  = {4},
  pages   = {515--534},
  year    = {1982},
  month   = dec,
  doi     = {10.1007/BF01457454}
}
}

\bib{
@inproceedings{10.1145/380752.380857,
author = {Ajtai, Mikl{\'o}s and Kumar, Ravi and Sivakumar, D.},
title = {A sieve algorithm for the shortest lattice vector problem},
year = {2001},
isbn = {1581133499},
publisher = {Association for Computing Machinery},
address = {New York, NY, USA},
url = {https://doi.org/10.1145/380752.380857},
doi = {10.1145/380752.380857},
booktitle = {Proceedings of the Thirty-Third Annual ACM Symposium on Theory of Computing},
pages = {601–610},
numpages = {10},
location = {Hersonissos, Greece},
series = {STOC '01}
}
}

\bib{
@inproceedings{10.1145/2746539.2746606,
author = {Aggarwal, Divesh and Dadush, Daniel and Regev, Oded and Stephens-Davidowitz, Noah},
title = {Solving the Shortest Vector Problem in $2^n$ Time Using Discrete Gaussian Sampling: Extended Abstract},
year = {2015},
isbn = {9781450335362},
publisher = {Association for Computing Machinery},
address = {New York, NY, USA},
url = {https://doi.org/10.1145/2746539.2746606},
doi = {10.1145/2746539.2746606},
booktitle = {Proceedings of the Forty-Seventh Annual ACM Symposium on Theory of Computing},
pages = {733–742},
numpages = {10},
keywords = {discrete Gaussian, lattices, shortest vector problem},
location = {Portland, Oregon, USA},
series = {STOC '15}
}
}

\bib{
@article{10.1016/S0020-0190(99)00083-6,
author = {Goldreich, O. and Micciancio, D. and Safra, S. and Seifert, J. -P.},
title = {Approximating shortest lattice vectors is not harder than approximating closest lattice vectors},
year = {1999},
issue_date = {July 30, 1999},
publisher = {Elsevier North-Holland, Inc.},
address = {USA},
volume = {71},
number = {2},
issn = {0020-0190},
url = {https://doi.org/10.1016/S0020-0190(99)00083-6},
doi = {10.1016/S0020-0190(99)00083-6},
journal = {Inf. Process. Lett.},
month = jul,
pages = {55–61},
numpages = {7},
keywords = {computational complexity, computational problems in integer lattices, linear error-correcting codes, reducibility among approximation problems, theory of computation}
}
}

\bib{
@book{MicciancioGoldwasser2002,
  author    = {Micciancio, Daniele and Goldwasser, Shafi},
  title     = {Complexity of Lattice Problems: A Cryptographic Perspective},
  publisher = {Kluwer Academic Publishers},
  series    = {The Kluwer International Series in Engineering and Computer Science},
  volume    = {671},
  year      = {2002},
  address   = {Boston, MA},
  isbn      = {978-0-7923-7688-0}
}
}

\bib{
@book{van1981another, title={Another NP-complete partition problem and the complexity of computing short vectors in a lattice}, author={van Emde-Boas, P.}, series={Report. Department of Mathematics. University of Amsterdam},  year={1981}, publisher={Department, Univ.} }
}

\bib{
@inproceedings{10.1145/1806689.1806739,
author = {Micciancio, Daniele and Voulgaris, Panagiotis},
title = {A deterministic single exponential time algorithm for most lattice problems based on Voronoi cell computations},
year = {2010},
isbn = {9781450300506},
publisher = {Association for Computing Machinery},
address = {New York, NY, USA},
url = {https://doi.org/10.1145/1806689.1806739},
doi = {10.1145/1806689.1806739},
booktitle = {Proceedings of the Forty-Second ACM Symposium on Theory of Computing},
pages = {351–358},
numpages = {8},
keywords = {lattice algorithms, Voronoi cell, SVP, SIVP, CVP},
location = {Cambridge, Massachusetts, USA},
series = {STOC '10}
}
}

\end{document}
